\section{本論文で示したこと}
\label{sec:conclusion-results}

本論文では，HHLアルゴリズムの数学的な構成を理解することを目的として，
量子計算の基礎から量子線形方程式の解状態までを順に説明した．

第2章では，有限次元複素内積空間の単位ベクトルとして量子状態を定義し，
ユニタリ変換，射影測定，複合系の公理を述べた．
テンソル積については，普遍性による定義から積基底を構成し，
線形写像の積，随伴およびユニタリ性に関する性質を証明した．
また，値レジスタに応じた制御回転をブロック対角なユニタリ演算として定義し，
補助量子ビットの測定確率と条件付き状態を導いた．

第3章では，量子フーリエ変換，量子位相推定，振幅増幅を説明した．
さらに，エルミート行列$A$に対する$e^{\mathrm{i}At}$がユニタリであり，
$A$の固有値$\lambda_j$が位相$\lambda_jt/(2\pi)$として現れることを証明した．
これにより，HHLで固有値を量子レジスタへ符号化するための準備が得られた．

第4章では，量子線形方程式問題の出力を
\[
 \ket x=\frac{A^{-1}\ket b}{\lVert A^{-1}\ket b\rVert}
\]
と定義した．位相推定，固有値の逆数に比例する制御回転，逆位相推定，
補助量子ビットの測定という各段階の状態を追い，
理想的な場合に成功時の出力が上の解状態となることを証明した．
また，$2\times2$行列の具体例では，固有値分解と位相を計算し，
成功確率が$5/8$，条件付き出力が
$(3\ket0-\ket1)/\sqrt{10}$となることを確認した．

以上により，HHLアルゴリズムにおける逆数変換は，
逆行列を行列として構成する操作ではなく，固有値の逆数を
量子状態の確率振幅へ反映する操作であることが分かる．

